% =========================================================
% Limits at Infinity — Notes
% Chapter: continuity
% Section: limits-at-infinity
% Volume III · Analysis
% =========================================================

\section{Limits at Infinity}
\label{sec:limits-at-infinity}

% ---------------------------------------------------------
% Toolkit
% ---------------------------------------------------------

\begin{tcolorbox}[title={Toolkit --- Limits at Infinity}]
\begin{itemize}
  \item $+\infty$ is adherent to $X$ iff $X$ is unbounded above
    ($\sup X = +\infty$).
  \item $\lim_{x \to +\infty} f(x) = L$: replace $|x - c| < \delta$
    with $x > M$. Same $\varepsilon$ condition on outputs.
  \item All standard limit laws carry over unchanged.
  \item Sequential characterisation: $\lim_{x \to +\infty} f(x) = L$
    iff $f(x_n) \to L$ for every sequence $x_n \to +\infty$ in $X$.
\end{itemize}
\end{tcolorbox}

% ---------------------------------------------------------
% Definition — Infinite Adherent Points
% ---------------------------------------------------------

\begin{tcolorbox}[title={Definition --- Infinite Adherent Points}]
\begin{definition}[Infinite Adherent Points]
\label{def:infinite-adherent-points}
Let $X \subseteq \mathbb{R}$.
\begin{itemize}
  \item $+\infty$ is \textbf{adherent to $X$} if
    $\forall M \in \mathbb{R},\; \exists x \in X : x > M$,
    equivalently $\sup X = +\infty$.
  \item $-\infty$ is \textbf{adherent to $X$} if
    $\forall M \in \mathbb{R},\; \exists x \in X : x < M$,
    equivalently $\inf X = -\infty$.
\end{itemize}
A set $X$ is bounded if and only if neither $+\infty$ nor $-\infty$
is adherent to $X$.
\end{definition}
\end{tcolorbox}

\begin{remark*}[Standard quantified statement]
\[
  +\infty \text{ adherent to } X
  \;\iff\;
  \forall M \in \mathbb{R},\;
  \exists x \in X : x > M
  \;\iff\;
  \sup X = +\infty.
\]
\end{remark*}

\begin{remark*}[Interpretation]
Infinite adherence is the unboundedness analogue of the finite
adherence condition. A finite point $x_0$ is adherent to $X$ when
$X$ has elements arbitrarily close to $x_0$. The point $+\infty$
is adherent to $X$ when $X$ has elements arbitrarily large. Both
are instances of adherence in the extended real line
$\overline{\mathbb{R}} = \mathbb{R} \cup \{-\infty, +\infty\}$,
where the neighbourhoods of $+\infty$ are the rays $(M, +\infty)$.
\end{remark*}

% ---------------------------------------------------------
% Definition — Limits at Infinity
% ---------------------------------------------------------

\begin{tcolorbox}[title={Definition --- Limit at Infinity}]
\begin{definition}[Limit at Infinity]
\label{def:limit-at-infinity}
Let $X \subseteq \mathbb{R}$ with $+\infty$ adherent to $X$, and
let $f : X \to \mathbb{R}$. The function $f$ \textbf{converges to
$L$ as $x \to +\infty$}, written $\lim_{x \to +\infty} f(x) = L$,
if for every $\varepsilon > 0$ there exists $M \in \mathbb{R}$ such
that
\[
  x \in X \;\wedge\; x > M
  \;\Rightarrow\;
  |f(x) - L| < \varepsilon.
\]
Similarly, $\lim_{x \to -\infty} f(x) = L$ if for every
$\varepsilon > 0$ there exists $M \in \mathbb{R}$ such that
\[
  x \in X \;\wedge\; x < M
  \;\Rightarrow\;
  |f(x) - L| < \varepsilon.
\]
\end{definition}
\end{tcolorbox}

\begin{remark*}[Standard quantified statement]
\[
  \lim_{x \to +\infty} f(x) = L
  \;\iff\;
  \forall \varepsilon > 0,\;
  \exists M \in \mathbb{R},\;
  \forall x \in X:\;
  x > M \Rightarrow |f(x) - L| < \varepsilon.
\]
\end{remark*}

\begin{remark*}[Negated quantified statement]
\[
  \lim_{x \to +\infty} f(x) \neq L
  \;\iff\;
  \exists \varepsilon_0 > 0,\;
  \forall M \in \mathbb{R},\;
  \exists x \in X:\;
  x > M \;\wedge\; |f(x) - L| \geq \varepsilon_0.
\]
\end{remark*}

\begin{remark*}[Comparison with finite limits]
The finite limit $\lim_{x \to c} f(x) = L$ uses the proximity
condition $|x - c| < \delta$. The limit at $+\infty$ replaces this
with $x > M$: instead of $x$ lying within $\delta$ of $c$, $x$
must exceed the threshold $M$. The output condition $|f(x) - L|
< \varepsilon$ is identical. All standard limit laws (sum, product,
quotient, composition, squeeze) carry over with the same proofs,
replacing $|x - c| < \delta$ by $x > M$ throughout.
\end{remark*}

\begin{remark*}[Interpretation]
A limit at infinity formalizes the intuition that $f(x)$ stabilizes
near $L$ as $x$ grows without bound. The threshold $M$ plays the
role of $c - \delta$ in the finite limit: everything to the right
of $M$ in $X$ must be within $\varepsilon$ of $L$. The definition
fits naturally into the framework of limits as $x_0$ approaches a
point adherent to the domain, here the adherent point being
$+\infty$.
\end{remark*}
